\section{Mechanism families}\label{sec:mechanisms}


\regime{Motion}{how deployed systems write, adapt, consolidate, and share state in practice.}
Mechanisms matter when they cover lifecycle operations. The literature on always-on agents has produced a large toolbox: memory managers that page personal history in and out of context, graph stores that index it for multi-hop recall, reflection systems that distill trajectories into lessons, skill libraries that turn successful executions into callable procedures, consolidation schedules that decide what to keep, and shared substrates that let agents read and write each other's state. Product by product, this looks like steady progress. Against the persistent-state lifecycle, the pattern is more uneven. Many mechanisms concentrate their strengths on the same three operations, write, organize, and retrieve, and leave the same governance fields weak: authority, scope, provenance, recoverability, and the validate/audit/rollback stages that enforce them. The mechanism families are therefore read by the operation each makes explicit, the governance field it leaves implicit, and the concrete consequence of that omission. Memory managers, reflection systems, skill induction, continual adaptation, and shared-memory systems together lead to the falsifiable controlled-compounding criterion at the end of the part.

Two pieces of prior scaffolding matter here. The cognitive-architecture lineage treats an agent as memory plus action plus a decision process, and persistent state extends that loop across sessions \citep{sumers2024cognitivearchitectureslanguageagents,yao2023reactsynergizingreasoningacting}, with the episodic, semantic, and procedural state types echoing the long-standing memory-literature distinction \citep{shinn2023reflexionlanguageagentsverbal}. What modern systems add is mechanism: explicit code for writing, indexing, consolidating, and reusing that state. What almost all of them still lack is governance: code for deciding whether a write is authorized, whether a lesson is true, whether a skill is still safe to call, and how to undo any of these once they prove wrong. Table~\ref{tab:mechanisms} makes this diagnostic concrete by mapping representative families to the operation each makes explicit and the governance field each leaves weak.

\begin{table}[t]
\centering
\small
\renewcommand{\arraystretch}{1.2}
\setlength{\tabcolsep}{4pt}\renewcommand{\arraystretch}{1.2}
\begin{tabularx}{\linewidth}{p{0.17\linewidth}YYY}
\toprule
Family & Operation made explicit & Governance field left implicit & Concrete consequence \\
\midrule
Memory managers and graph memory & Write, organize, retrieve, consolidate; graph stores add multi-hop and temporal indexing & Authority and recoverability: who authorized a fact to act, can deletion be audited & Deletion edits the primary store while summaries, embeddings, and promoted tiers persist \\
Reflection and experience learners & Update: convert traces into lessons or causal abstractions & Validation: is the induced lesson true & One failed trajectory becomes a plausible but unproven durable rule \\
Skill and workflow induction & Procedural write and reuse of executable skills & Authority and mutability: test, scope, deprecate before reuse & A trajectory becomes a callable action whose misuse has tool consequences \\
Continual adaptation & Plasticity at write and update across task streams & Stability: validate, audit, forget, rollback & Helpful adaptation turns net-negative; competence and disorder accumulate together \\
Multi-agent and shared memory & Shared state, schema-bound mutation, access control & Scope and consistency: revocation and propagation & Corrupt or stale state crosses users or agents through shared memory \\
\bottomrule
\end{tabularx}
\caption{Mechanism families through the lifecycle lens. Each family makes one operation explicit and leaves a governance field (authority, scope, mutability, provenance, recoverability) implicit, with a concrete consequence. The comparison identifies which state-governance gaps recur across the surveyed works rather than ranking memory products.}
\label{tab:mechanisms}
\end{table}

\subsection{Memory managers and graph memory}

The first and largest family treats memory as something to be paged, indexed, and retrieved. Operating-system-inspired managers make the write, organize, and retrieve operations explicit and efficient: they externalize long-term state, manage a bounded working set, and decide what moves between tiers \citep{packer2024memgptllmsoperatingsystems,kang2025memoryosaiagent}. Agentic note systems push organization further by letting stored entries link and refactor each other as new information arrives, so that structure is emergent rather than fixed at write time \citep{xu2025amemagenticmemoryllm}. Graph-backed stores add structured multi-hop and temporal indexing on top of raw recall: knowledge-graph memory with neurobiologically-inspired retrieval supplies association structure that flat vector stores lack \citep{gutierrez2024hipporag}, and temporal knowledge graphs attach explicit validity windows so that facts can be queried as of a point in time \citep{rasmussen2025zep}. These systems are genuinely effective at what they make explicit: they scale personal history beyond the context window, recover semantically distant facts, and support reasoning that chains several stored items.

What this family leaves implicit is authority and recoverability. A managed or graph-indexed store records that a fact exists and how to retrieve it, but rarely records who authorized that fact to influence an action, under what scope, or how to prove it was later removed. The concrete consequence is a deletion gap that the lifecycle frame predicts: when a user or policy demands erasure, the operation edits the primary store, while the summaries, embeddings, promoted semantic tiers, and cached retrieval traces derived from the deleted item persist and can still be retrieved. Graph structure makes the recoverability problem more acute, because a single deleted node may have already propagated through edges into linked notes and consolidated summaries, and the typical graph store offers no cascade that follows those edges. The deeper representational failure is that flat-text long-term memory loses track of which source a fact came from and how it should be used, a source-monitoring breakdown that one analysis names provenance-role collapse and shows is mitigated only by typed memory representations rather than by better retrieval \citep{jin2026provenancerolecollapse}. The gap this family exposes for the thesis is direct: a memory manager optimizes accumulation and retrieval, the plasticity-side operations, while the governance fields that make state safe to act on, authority over a write and an auditable path to undo it, are left to the surrounding system that usually does not implement them. Retrieval that works is not the same as state that is governed.

\subsection{Reflection and experience learners}

The second family makes the update operation explicit. Rather than storing raw traces, these systems convert experience into reusable abstractions: verbal self-reflections that diagnose what went wrong on a trial \citep{shinn2023reflexionlanguageagentsverbal}, natural-language insights mined across many past episodes \citep{zhao2024expelllmagentsexperiential}, and causal abstractions that a continually learning agent refines so that later episodes inherit what earlier ones discovered \citep{majumder2023clincontinuallylearninglanguage}. A complementary line replays past trajectories directly as in-context experience rather than distilling them, integrating raw episodes into the working context at decision time \citep{liu2025contextualexperiencereplay}, and reasoning-memory systems generalize this by storing solution patterns that scale self-improvement on new tasks \citep{ouyang2025reasoningbank}. The shared move is attractive because it is the most direct route to compounding competence: an agent that turns each outcome into a lesson should, in principle, get monotonically better.

The governance field this family leaves implicit is validation. A reflection, insight, or causal rule is written into durable state on the strength of the agent's own judgment about a trajectory, with no separate gate that asks whether the induced lesson is actually true. The concrete consequence is that a single failed or misread trajectory becomes a plausible but unproven durable rule, and because the rule reads as a confident generalization, later retrieval surfaces it as if it were established fact. This is not hypothetical: the self-confirmation trap describes exactly the case where an agent that judges its own trajectories writes self-consistent-but-wrong experiences into memory, and the proposed remedy decouples writing from judging through heterogeneous executors, a third-party distiller, and a consensus gate that decides which experiences are committed \citep{zhu2026selfconfirmation}. The lesson generalizes the lifecycle claim: reflection systems are strong on update and weak on validate, so they convert plasticity into durable belief without the stability check that would keep a false lesson from compounding. The gap is that few of these systems, among those we surveyed, carry the provenance needed to later identify which reflection caused a downstream regression, or the rollback handle needed to revert it; an advisory lesson that degrades gracefully is forgiving of this omission, but as the next subsection shows, the same omission is far less forgiving once the durable artifact is executable.

\subsection{Skill and workflow induction}

The third family makes procedural write and reuse explicit: it turns a successful trajectory into a callable skill or a reusable workflow abstraction. The canonical embodied example grows a library of executable skills over an open-ended lifetime, so that later tasks call procedures discovered earlier \citep{wang2023voyageropenendedembodiedagent}; the canonical web example induces high-level workflow abstractions from past trajectories and applies them in new sessions \citep{wang2024agentworkflowmemory}. Recent work treats the skill store as a first-class artifact with its own engineering discipline. Procedural-memory systems add deprecation and update support so that stored workflows are not write-once \citep{fang2026mempexploringagentprocedural}; skill-management frameworks model the library as a self-maintaining software ecosystem with typed contracts to control accumulating technical debt \citep{pu2026skillops}; and learned-skill methods add a verification gate that checks a candidate skill before it is written and a score-based maintenance pass that keeps the library compact \citep{mi2026skillpro}. Selective-formalization approaches go further toward durability by recording reusable steps as auditable, versioned notebooks with explicit validation gates and fallback paths \citep{elhattami2026skillnb}.

This family exposes the actionability and authority axes most sharply, because a procedural skill is not advisory: it is an executable commitment whose misuse has tool consequences. The governance fields left implicit are authority and mutability, the questions of whether a skill is still authorized to run and whether it has been tested and scoped against current conditions before reuse. The concrete failure mode is skill drift, where a skill that was correct when written silently becomes wrong as the environment changes; one analysis formalizes this as a contract violation and validates the role-bearing environment assumptions of a stored skill against live conditions to trigger proactive repair \citep{fan2026skilldriftcontract}. The danger is amplified because a static skill can look benign and only turn harmful in the presence of accumulated persistent state, which is why runtime auditing that probes stored skills under targeted live conditions catches failures that static review misses \citep{lan2026runtimeskillaudit}. Whether a refined skill even transfers, rather than quietly overfitting to the role it was learned in, is itself an open question that a dedicated benchmark isolates, finding that procedural skills often become role-specialized and lose effectiveness off-distribution \citep{belikova2026managingproceduralmemory}, and a separate probe shows that embedding-based procedural retrieval discards the temporal ordering structure that procedures depend on, producing a generalization cliff \citep{kohar2025proceduralmemoryretrieval}. The comparison is instructive: most skill libraries optimize cheap reuse and do not re-test a stored procedure before calling it, so among representative systems the ones closest to lifecycle-complete are precisely those that add the missing stages, MemP with explicit deprecation, SkillOps with typed contracts and maintenance, Skill-Pro with a write-time verification gate, and SkillNB with versioned validation. The gap this family exposes is the steepest in the survey: the actionability of a procedural write raises the authorization, validation, and rollback bar exactly where mechanisms are weakest, so an ungoverned skill library accumulates executable commitments faster than it accumulates any way to verify, scope, or revoke them.

\subsection{Continual adaptation and the plasticity-stability problem}

The first three families each name a way to compound competence. Continual adaptation names why compounding is hard, and the harm it causes when ungoverned. The tension is the classic stability-plasticity tradeoff of continual learning \citep{wang2023continuallearningsurvey}, and the interference it produces traces back to the catastrophic-interference result for sequential learning in connectionist networks \citep{mccloskey1989catastrophic}: weights tuned for new tasks overwrite the representations that solved old ones. Surveys of lifelong learning for LLMs map the response space, in-context adaptation, parameter update, and memory augmentation, across this tradespace \citep{zheng2024lifelongllmsurvey}, and surveys of self-evolving agents catalogue the loops that chain exploration, skill extraction, and internalization \citep{gao2025selfevolvingagentssurvey}. Always-on agents face the same tradeoff in an externalized form. Where lifelong-learning surveys catalogue internal weight updates, these agents adapt mostly through retrieval and tool-grounded state, using the reflections, experience buffers, and skill libraries of the previous subsections rather than gradient steps.

\subsubsection{Parametric versus non-parametric adaptation}

The parametric alternatives sharpen the contrast and clarify why the externalized route is attractive. Test-time training updates model weights at inference, either from retrieved nearest neighbors \citep{hardt2023tttnn} or from the few-shot examples in the current task \citep{akyurek2024ttt}, boosting generalization without a permanent commitment. Knowledge editing rewrites facts in place, but editing at scale itself induces gradual and catastrophic forgetting \citep{gupta2024editingforgetting}: the same plasticity-stability failure reappears in weight space, now without the provenance or rollback handles that externalized state can keep, since a weight edit leaves no record of which fact it overwrote or how to restore it. This is the core trade. Parametric memory is efficient and fast to consult but fragile to edits and opaque to audit; non-parametric external memory is costlier per query but inspectable and, in principle, reversible. A longitudinal study that distinguishes parametric memory depth from retrieval memory access uses surprise- and valence-gated consolidation to test whether goal-conditioned tendencies persist after the working context is unloaded \citep{han2026memorydepth}, and an empirical lifelong benchmark finds that nonparametric memory outperforms parametric weight updates while catastrophic forgetting grows with interaction length \citep{fan2025lifestatebench}. Always-on agents take the externalized route precisely to make adaptation inspectable and reversible; the lifecycle frame then makes the tension exact. Plasticity lives at write and update; stability is enforced at validate, audit, and forget; an agent that exposes the first set of operations without the second accumulates risk even as its task scores rise. We state the criterion in falsifiable form: adaptation is net-positive only when the system can identify, de-authorize, and revert the specific state update that later caused a regression. A system that cannot perform that revert has not shown controlled compounding, however high its aggregate score climbs.

\subsubsection{How helpful adaptation turns net-negative}

Several reported patterns show the criterion being violated in practice. \emph{Skill shadowing}: a procedural skill that succeeded early is reused after the environment or goal has shifted, so a correct past execution becomes a wrong current one, and a library built for cheap reuse does not re-test a stored procedure before calling it \citep{wang2023voyageropenendedembodiedagent}, the drift the skill subsection formalized. \emph{Interference}: writes that help recent tasks degrade retention of older ones, the backward-transfer failure of continual learning expressed here through retrieval rather than weights. \emph{Non-monotonic self-evolution}: iterating an explore-extract-internalize loop does not guarantee improvement \citep{gao2025selfevolvingagentssurvey}, and frozen-weight agents degrade over time through distinct aging mechanisms, compression, interference, revision, and maintenance, that a longitudinal benchmark isolates and targets with stage-specific repair \citep{zhu2026aging}. \emph{Replay weakness}: replay is the most direct compounding mechanism \citep{liu2025contextualexperiencereplay}, yet a lifelong agent benchmark finds naive replay weak for LLM agents because irrelevant retrieved episodes consume context, with structured aggregation helping more than raw recall \citep{zheng2025lifelongagentbenchevaluatingllmagents}, and an environment-drift benchmark that frames the world as a sequence of progressive updates finds that agents average well under forty percent at keeping memory aligned to a shifting world \citep{xu2026evoarena}.

\subsubsection{The baseline-beats-memory finding}

Recent controlled studies sharpen these patterns into a baseline question that is uncomfortable for the whole field. A continual-learning benchmark built on a learnable latent structure that only stateful systems can exploit finds that dedicated memory systems do not reliably improve learning, and that naive in-context use of recent history outperforms several purpose-built memory architectures \citep{asawa2026clbench}. A separate benchmark with controlled reusable task streams shows that memory-induced degradation is hidden by naive evaluation streams and surfaces only when reuse is controlled, and pairs the finding with a filter for unreliable experience \citep{shu2026agentcl}; a long-horizon benchmark over continuous agent-environment interaction across six domains tells the same story \citep{zhao2026amabench}, and a co-play-derived task suite that instruments memory reads, writes, and validators exposes where agents fail to persist memory across tasks at all \citep{doss2026minenpc}. The starkest result is that continuously LLM-consolidated textual memory has a utility curve that rises and then degrades below the no-memory baseline as consolidation accumulates \citep{zhang2026usefulmemoriesfaulty}: the memory stops helping and then actively harms, and the harm is a function of how much has been written, not of any single bad write. This is the silent-entropy reading of long-run degradation, accumulating disorder that grows with interaction count and is invisible to any metric that scores only the current answer. The collective message is empirical: stored state helps only when it is governed, not when it is simply accumulated, and an ungoverned memory system can be strictly worse than having no memory at all once the horizon is long enough.

\subsubsection{Forgetting and consolidation as the stability operators}

Principled forgetting and consolidation are the stability operators that should contain harmful plasticity, and a wave of recent work develops those operators. The shared insight is that forgetting should be governed, not incidental. Several systems model forgetting as decay in accessibility rather than destructive deletion: a dual-layer hierarchy applies differential exponential-decay rates keyed to relevance, access frequency, and temporal pattern \citep{wei2026fademem}, while a read-path-versus-control-path split gates recall by agent uncertainty and treats forgetting as reduced accessibility \citep{rana2026oblivion}. A selective-forgetting taxonomy organizes the design space into passive decay, active deletion, safety-triggered removal, and adaptive reinforcement, making explicit that different mechanisms preserve and discard different knowledge \citep{gu2026fsfm}, and a retention-schema approach grounds six forgetting policies in theory and evaluates them on a dedicated benchmark \citep{alqithami2025mars}. Consolidation work decouples fast acquisition from slow integration: an RL-trained offline consolidator separates per-session memory capture from cross-session consolidation and shrinks the active bank while transferring across environments \citep{ye2026autodreamer}, a human-inspired design combines sleep-phase consolidation, interference-based forgetting, engram maturation, and reconsolidation on retrieval \citep{kerestecioglu2026humaninspiredmemory}, and a predictability framing reframes what to retain as a prediction-error problem learned from interaction rather than a fixed heuristic \citep{ma2026nemori}. Other systems make consolidation governable through structure: a short-to-long-term consolidation policy frames retention as a per-triple keep-or-drop Q-learning decision over a temporal knowledge graph \citep{kim2026stltransfer}, residual-tree storage lets related experiences share a base while keeping incremental deltas \citep{tan2026deltamem}, deferred extraction waits for sustained recurrence before paying for expensive consolidation \citep{dai2026recmem}, and a control-plane analysis shows that where the LLM sits, on the recall plane versus a control plane that can supersede, release, and purge, determines which forgetting failure mode dominates \citep{yang2026controlplaneforgetting}. A role-specific memory policy even learns from periodic QA feedback and retains a failure-rollback path so that a bad consolidation can be reverted \citep{chen2026adamem}, which is the rare system that closes the loop the falsifiable criterion demands.

These findings reframe continual adaptation as a state-governance problem, more than a learning-rate problem. The interventions that keep compounding controlled are precisely the underused lifecycle stages: validation gates before a lesson becomes a durable rule, provenance and recency tags so that superseded skills and facts lose authority rather than lingering, and audit and rollback so that a degrading update can be identified and reverted. The gap this subsection exposes is that the works we coded contain rich plasticity mechanisms and an increasingly rich vocabulary of forgetting, but still couple them weakly to authority and rollback: many forgetting systems decide what to decay without recording who authorized the decay or how to undo it if the decay was wrong, and benchmarks that interleave dependent sessions surface the degradation without scoring the write, deprecate, and rollback decisions that determine whether adaptation compounds or corrodes. A complementary line tries to absorb non-parametric experience back into parametric memory under control, co-updating an experiential rule pool and the policy weights from the same trajectories so that stable behaviors migrate from the rule pool into the model \citep{ye2026jerp}, and model-based continual RL governs a dual recency-plus-diversity replay buffer under a fixed memory budget to bound interference \citep{alyahya2026arrow}; both are promising precisely because they make the migration auditable and budget-bounded rather than leaving it to accumulate unchecked.

\subsection{Multi-agent and shared-memory mechanisms}

The final family moves from a single agent's lifecycle to state shared across many. Multi-agent frameworks have long used a shared substrate as de-facto memory: a publish-subscribe message pool acts as untyped shared state in canonical SOP-driven coordination \citep{hong2023metagpt}, and generative-agent sandboxes showed that observation-reflection-retrieval memory in a shared environment produces believable social behavior \citep{park2023generativeagentsinteractivesimulacra}. Interactive social-evaluation frameworks made shared world state and relationship history a core measurement target \citep{zhou2024sotopiainteractiveevaluationsocial}. The newer wave makes the operation this family makes explicit, shared write with access control, into a first-class design concern. Access-governed shared memory enforces explicit authority and scope on multi-agent reads and writes \citep{margalit2026governedsharedmemorymultiagent}; multi-user collaborative memory adds dynamic access-control gates with field-level and role-anchored granularity \citep{rezazadeh2025collaborativememorymultiusermemory}; schema-grounded mutation enforces type-safe transitions with full audit trails to bridge untyped updates and reliable shared state \citep{zhang2026patchboardschemagroundedstatemutation}; and role-aware context routing governs which subset of shared memory each agent sees under a token budget, a selective-view form of scope control at retrieval time \citep{liu2025rcrrouter}.

The governance fields this family leaves implicit are scope and consistency: specifically, revocation and propagation control. The concrete consequence is that corrupt or stale state crosses users or agents through the shared substrate, and the failure-mode literature documents the cost. Jailbreak and poison propagation spreads across agent-society topologies, and shared-memory architectures propagate it more readily than independent-memory ones \citep{men2024troublemaker}; evaluator bias in stored trajectories propagates cross-temporally to future agents sharing memory, with no safe contamination threshold even under oracle consolidation, a failure named memory contagion \citep{liu2026memorycontagion}; and a thousand-scenario benchmark quantifies privacy leakage through inter-agent messages, shared memory, and tool arguments \citep{elyagoubi2026agentleak}. These results show that per-agent safety checks are insufficient when the memory substrate itself can carry compromised state across the scope boundary, which is the P6 keys point: shared state without enforced scope non-expansion lets one agent's authority and one user's private facts reach contexts they were never licensed for.

A second strand addresses consistency and forgetting in the shared setting directly, and it is here that the field's informal notion of shared memory most needs the formal consistency vocabulary of distributed systems. When many agents read and write the same memory stores and tool registries concurrently, the substrate inherits the classical anomalies of replicated state: a recent machine-checked study makes this precise, exhibiting the first formal consistency hierarchy for multi-agent LLM runtimes and naming anomalies such as stale-generation and phantom-tool that arise when concurrent writes race over shared state \citep{khan2026concurrency}. The trade-off behind these anomalies is the familiar consistency-versus-availability tension, which a concurrent-abstract-state-machine formalization renders as an explicit framework for deciding how much consistency a given replication policy actually needs, a decision a governed agent memory must make rather than leave implicit \citep{schewe2019replicated}. Database concurrency control supplies the matching proof technique: serializability is the standard correctness criterion for concurrent access, and a machine-checked CV-rules framework for verifying it can be adapted to certify that interleaved multi-agent memory mutations remain equivalent to some safe serial order \citep{hoshino2026cv-rules}. These formal models give the scope-and-consistency gap a precise statement: an ungoverned shared substrate offers neither a defined consistency level nor a serializability guarantee, so concurrent writes from different agents can interleave into states that no single agent ever intended. Shared state need not always mean contention: read-heavy multi-agent workloads can share compressed key-value cache state with minimal coherency overhead, cutting memory roughly ninety-eight percent relative to isolated caches, which shows scope-bounded sharing and efficiency are compatible when the access pattern is governed \citep{patel2026polykv}. Where some agents may be faulty or adversarial rather than only concurrent, the requirement strengthens to fault-tolerant agreement: a Byzantine-consensus protocol that races on protocol structure rather than timeout-based delays shows that agents can still agree on shared state under adversarial participants \citep{giridharan2026ambulance}, and a diagnostic layer that combines policy analysis with Bayesian trust can flag a Byzantine agent before it corrupts the shared channel, supplying the detection half of a consistency guarantee \citep{murimi2026bard}. A co-forgetting protocol combines semantic-relevance voting and multi-scale temporal decay ratified by Byzantine consensus so that shared-memory pruning is safe under disagreement \citep{bach2025coforgetting}, and a phylogenetic version-control scheme coordinates decentralized agents through an asynchronous shared repository with humans promoting and pruning branches as a selection signal \citep{huang2025evogit}. Decentralized stigmergic coordination shows that agents can cohere above a critical density by depositing environmental traces without any central controller, but the same work notes the substrate is open to pollution and lacks audit mechanisms to verify trace authenticity or recover from coordinated poisoning \citep{khushiyant2025collectivememory}. The social and organizational dimension adds its own state: multi-party conversational memory requires speaker attribution and turn-by-turn tracking that single-user memory can ignore \citep{yang2026groupmembenchbenchmarkingllmagent}, and at the longest horizons, continuously running multi-agent platforms surface governance divergence and outright collapse under identical setups over weeks to months \citep{akkil2026emergenceworld}. Two framings try to make this institutional state governable: recasting agent runs as a learning organization whose experience documents inherit across runs under safeguards against knowledge degradation \citep{xie2026forage}, and reframing shared-memory governance as artificial selection over candidate state, persist, keep private, reject, revise, or supersede, with evaluation axes for correction pathways and provenance fidelity \citep{cuadros2026governedselection}.

The gap this family exposes is the multi-agent specialization of the survey's core claim. Shared substrates make write and access-control explicit, and a few make schema-bound mutation and consensus pruning explicit, but revocation, deletion propagation across agents, and rollback traceability remain thin: when one speaker leaves a group or one delegation is revoked, the dependent state they wrote into the shared substrate, and everything consolidated from it, typically has no cascade to follow and no audit trail to reconstruct. Corrupt state therefore propagates faster than any mechanism to contain it, which is why long-horizon multi-agent platforms measure collapse empirically but cannot yet explain or repair it mechanistically. The more governance-aware systems here, governed shared memory, schema-grounded mutation, and consensus-ratified co-forgetting, begin to instrument scope and consistency, but none yet closes the loop from a revoked authority back through every shared and derived copy of the state it licensed.

\subsection{The controlled-compounding criterion}

Read together, Table~\ref{tab:mechanisms} and the five families form a single diagnostic rather than five separate reviews. Many mechanisms share the same strengths, write, organize, and retrieve, and the same weaknesses, the validate, audit, and rollback stages that would govern what those operations produce. Memory managers leave the authority of a written item implicit, so deletion edits the store instead of producing auditable erasure. Reflection systems make update explicit but leave validation implicit, so one failed trajectory becomes a durable false rule. Skill libraries expose the actionability axis most sharply, turning trajectories into executable commitments while leaving authorization and deprecation weakest. Continual adaptation shows that without the stability operators the act of accumulating helps until it harms. Shared substrates carry all of these omissions across the scope boundary at once. The response implied by the table is \emph{controlled compounding}: accepting a write or update only when it preserves the five invariants, so that competence accumulates only on state that is grounded, scoped, testable, and reversible. The criterion is deliberately falsifiable. A system claims controlled compounding only if, when a later regression is traced to an earlier write, reflection, skill, consolidation, or shared-memory entry, that specific update can be identified, de-authorized, and reverted with its derived state. By that standard most current mechanisms have not demonstrated controlled compounding; they have demonstrated accumulation. The difference between the two is the governance this survey argues the field should measure.
